% Stay_v2_algorithm_blended.tex
% Neurema — Adaptive Revision Scheduling (v2) — Blended & cleaned LaTeX source
% Merges Stay_v1 parameter math with the hybrid bubble scheduler from Stay_Effective_final.
% Author: Neurema Scheduler Assistant
% Version: v2.1 — blended (2025-11-08)

\documentclass[11pt]{article}
\usepackage{amsmath}
\usepackage{booktabs}
\usepackage{geometry}
\usepackage{enumitem}
\geometry{margin=1in}
\title{Neurema — Adaptive Revision Scheduling (v2.1)\\LaTeX Specification}
\author{Generated by Neurema Scheduler Assistant}
\date{2025-11-08}

\begin{document}
\maketitle
\begin{abstract}
This document specifies the blended v2 scheduler: it merges the Stay\_v1 parameter definitions and math (EF, PI, CRS) with the capacity-capped, formulaic bubble scheduler used in the 9-month "Stay\_Effective" runs. The spec contains exact formulas, worked numeric examples, clear variable names, capacity knobs, and concise, implementable pseudocode.
\end{abstract}

\section{Notation and conventions}
\begin{description}[leftmargin=!,labelwidth=3cm]
  \item[Indices] Scalars are upper-case (EF, PI). Time is measured in days. Probabilities are in $[0,1]$.
  \item[Normalization] Raw inputs map to $[0.6,1.2]$ unless otherwise stated.
  \item[Constants] Adjustable system knobs are capitalised and documented in \S\ref{sec:params}.
\end{description}
\section{Core Parameters}
Each topic is evaluated using three core performance parameters:

\begin{itemize}
\item \( RT \): Time Ratio \( = \frac{T_{expected}}{T_{taken}} \)
\item \( D \): Difficulty Index (subjective or algorithmic difficulty rating)
\item \( AS \): Accuracy Score Index (based on accuracy percentage)
\end{itemize}

All indices are mapped between 0.6 and 1.2 for normalization.

\section{Efficiency Factor (EF)}
The Efficiency Factor measures micro-level learning performance for a topic.

\begin{equation}
EF = RT + D + AS
\end{equation}

\subsection*{(a) Time Ratio Index (RT)}
\begin{tabular}{ccc}
\toprule
Time Ratio (\(T_E/T_T\)) & Interpretation & Index (RT) \\
\midrule
\(\ge 1.2\) & Much faster than expected & 1.1 \\
\(1.0 - 1.2\) & Slightly faster & 1.0 \\
\(0.8 - 1.0\) & On time / slightly slower & 0.9 \\
\(0.6 - 0.8\) & Slower than expected & 0.8 \\
\(< 0.6\) & Very slow & 0.7 \\
\bottomrule
\end{tabular}

\subsection*{(b) Accuracy Score Index (AS)}
\begin{tabular}{ccc}
\toprule
Accuracy (\%) & Interpretation & Index (AS) \\
\midrule
\(\ge 90\%\) & Excellent mastery & 1.2 \\
80--89\% & High accuracy & 1.1 \\
70--79\% & Moderate understanding & 1.0 \\
60--69\% & Weak grasp & 0.9 \\
\(< 60\%\) & Poor performance & 0.8 \\
\bottomrule
\end{tabular}

\subsection*{(c) Difficulty Index (D)}
\begin{tabular}{ccc}
\toprule
Difficulty Level & Interpretation & Index (D) \\
\midrule
Easy & Low cognitive load & 1.2 \\
Moderate & Normal effort & 0.9 \\
Hard & Conceptually dense & 0.6 \\
\bottomrule
\end{tabular}

\subsection*{(d) EF Range Interpretation}
\begin{tabular}{ccc}
\toprule
EF Range & Topic Type & Meaning \\
\midrule
\(< 2.6\) & Struggling & Poor retention -- early revision \\
2.6--3.0 & Average & Normal grasp -- standard interval \\
\(> 3.0\) & Confident & Strong understanding -- revise later \\
\bottomrule
\end{tabular}
Each raw input is mapped to a normalized index in the range $[0.6,1.2]$. The Efficiency Factor is:
\[
EF = RT + D + AS
\]
Example: $RT=1.1,\ D=0.9,\ AS=1.1 \Rightarrow EF=3.1$.

\subsection{Possibility Index (PI)}
\[
PI = \log_{10}\left(1 + \frac{ND}{NS\cdot T_{min}}\right)
\]
Where:
\begin{itemize}
  \item $ND$ = number of days available until the target exam date.
  \item $NS$ = number of study slots (or sessions) per topic set.
  \item $T_{min}$ = minimum time fraction per session (days expressed as fraction of a day).
  
  Typical \( T_{min} \) values (in days per topic):
\begin{tabular}{ccc}
\toprule
Group & Description & \(T_{min}\) \\
\midrule
Major & Dense, heavy subjects & 0.10 \\
Medium & Moderate load subjects & 0.06 \\
Minor & Light, recall subjects & 0.03 \\
\bottomrule
\end{tabular}
\end{itemize}
Example: $ND=180,\ NS=50,\ T_{min}=0.10 \Rightarrow PI=\log_{10}(37)\approx1.568$.

\subsection{Composite Readiness Score (CRS)}
\[
CRS = \frac{EF\times PI}{\kappa}
\]
Where $\kappa$ is a scaling constant (default $\kappa=1.5$). Example: $CRS\approx3.24$.

\section{From CRS to effective interval}
We convert readiness into an effective interval using a modulation factor that depends on EF:
\[
\alpha_M(EF) = 0.5 + \frac{1}{1+e^{-\gamma (EF-\mu)}}
\]
Default sigmoid parameters: $\gamma=3,\ \mu=2.8$. Then:
\[
I_{eff} = BaseInterval \times CRS \times \alpha_M(EF)
\]
Example: $BaseInterval=5$ days, $EF=3.1 \Rightarrow \alpha_M\approx1.21,\ I_{eff}\approx19.6$ days.

\section{Mapping $I_{eff}$ to an exponential spacing model}
We use an exponential forgetting/spacing model:
\[
P(t)=e^{-t/S}\quad\Rightarrow\quad t=-S\ln(P^*)
\]
Choose $S$ so that $I_{eff}$ corresponds to a chosen reference probability $P_{ref}$ (default $P_{ref}=0.9$):
\[
S = -\frac{I_{eff}}{\ln P_{ref}}
\]
With $I_{eff}=19.6$ and $P_{ref}=0.9$, $S\approx186.1$ days. Then desired revisit times for other probability targets $P^*$ follow $t=-S\ln(P^*)$.

\section{Capacity-aware hybrid bubble model}
Scheduler must respect daily capacity constraints and support two revision types: single-topic (individual) and bubbles (grouped revisions). Bubbles reduce overhead by revising similar/hard topics together.

\subsection{System knobs and defaults}\label{sec:params}
\begin{itemize}
  \item $CAP_{topics/day}$: maximum individual topic events per day (default 20).
  \item $CAP_{bubbles/day}$: maximum bubbles per day (default 5).
  \item $Bubble_{size}^{min},Bubble_{size}^{max}$: min/max members per bubble (default 5--10).
  \item $BaseInterval$: baseline interval unit (days) used in $I_{eff}$ (default 5).
  \item $P_{4},P_{5}$: probability targets for successive spaced repetitions (e.g., 0.75 and 0.6). These are scheduling targets used to compute next desired days.
\end{itemize}

\subsection{High-level flow}
\begin{enumerate}
  \item For each topic compute normalized inputs $\rightarrow$ $EF, PI, CRS$ (Stay\_v1 math).
  \item Compute $\alpha_M(EF)$ and $I_{eff}$, then $S$ and desired revisit days for $P_4,P_5,\dots$.
  \item Insert each "next desired event" into a PendingPool as tuples $(day, priority=CRS, topicID)$.
  \item Daily scheduling (for day $d$):
    \begin{enumerate}
      \item Compute remaining individual topic slots: $slots = CAP_{topics/day} - individual\_events[d]$.
      \item Allow up to $CAP_{bubbles/day}$ bubbles on day $d$.
      \item Pop candidates from PendingPool with $desired\_day \le d$, ordered by priority (CRS) desc.
      \item While bubble slots $>0$ and slots $>0$:
        \begin{itemize}
          \item Form a bubble of up to $Bubble_{size}^{max}$ hardest remaining topics that are "alike" (same subject/section/skills). If not enough "alike" topics available, fill with highest-CRS leftovers.
          \item Schedule bubble on day $d$ (consumes 1 bubble slot and $m$ topic slots where $m$ is bubble size).
          \item For each member, compute next desired day using $t=-S\ln(P_{next})$ and re-insert into PendingPool.
        \end{itemize}
      \item After bubbles, schedule remaining highest-priority topics as individual events until daily topic slots are filled.
      \item Reinsert any un-scheduled candidates back into PendingPool (they will be reconsidered on later days).
    \end{enumerate}
\end{enumerate}

\section{Worked numeric example}
Using the running example from Stay\_v1:
\begin{itemize}
  \item $RT=1.1, D=0.9, AS=1.1 \Rightarrow EF=3.1$.
  \item $ND=180, NS=50, T_{min}=0.10 \Rightarrow PI=1.568$.
  \item $CRS = (3.1\times1.568)/1.5 \approx 3.2405$.
  \item $\alpha_M\approx 1.21,\ I_{eff}=BaseInterval\times CRS\times\alpha_M \approx 19.6$ days.
  \item $S = -I_{eff}/\ln 0.9 \approx 186.1$ days.
  \item If $P_4=0.75$, desired $t_4 = -S\ln(0.75) \approx 53.9$ days after last individual exposure.
\end{itemize}

\section{Implementation notes and caveats}
\begin{itemize}
  \item Topic "alikeness" (for bubble formation) must be defined by taxonomy: subject, chapter, skill type. Use semantic clustering or manual labels. 
  \item If daily capacity is tight, CRS ordering ensures higher-need topics are prioritized.
  \item Parameter tuning: $\kappa, BaseInterval, P_{ref}, CAP_{*}$ should be adjusted empirically on simulation data (recommend grid search / bayesian optimisation on historical runs).
  \item Edge cases: topics with extremely high CRS that continually push outothers should be allowed a fairness correction (e.g., max_wait cap) to avoid starvation.
\end{itemize}

\section{Version history}
\begin{tabular}{ll}
v2.0 & original blended spec (Stay\_v1 + Stay\_Effective) \\
v2.1 & cleaned duplicate blocks, added knobs, clarified bubble selection and fairness notes (this file)
\end{tabular}

\end{document}
